% Lecture Slides: a 16:9 Beamer deck for one lecture of a university course.
% The theme (colours, strip, blocks, code style, agenda) lives in
% beamerthemelecture.sty; this file holds only the lecture.
% For printable handouts, add `handout' to the class options: every step and
% the quiz answer then appear on a single page per slide.
\documentclass[aspectratio=169,11pt,t]{beamer}

\usetheme{lecture}
\usepackage{amsmath}
\hypersetup{hidelinks}

% Course and lecture data: the code and number appear on the title slide and
% in the footer of every slide.
\course{CS 241}{Data Structures and Algorithms}
\lecturenumber{7}
\title[Hash tables]{Hash tables}
\subtitle{Constant-time lookup, collisions and the load factor}
\author{Dr.\ Amara Nwosu}
\institute{School of Computing, Northfield University}
\date{Tuesday 14 October 2026}

% The parts of today's lecture, in order; the agenda and dividers use this list.
\lcsections{Hashing,Collisions,Load factor and resizing,Practice}

\begin{document}

% ------------------------------------------------------------------ title
{\lcstripfalse
\begin{frame}[plain,noframenumbering]
  \titlepage
\end{frame}}

\begin{frame}{Today}
  \begin{columns}[T,onlytextwidth]
    \column{0.48\textwidth}
      \vspace{2mm}
      \lcagenda
    \column{0.46\textwidth}
      {\color{lcMuted}\lclabel{Learning goals}}\par\vspace{2mm}
      {\small
      \begin{enumerate}
        \item Explain how a hash function maps keys to slots.
        \item Compare chaining with open addressing.
        \item Predict the cost of a lookup from the load factor.
        \item Decide when a table should grow.
      \end{enumerate}}
  \end{columns}
\end{frame}

% ------------------------------------------------------------------ part 1
\section{Hashing}

\begin{frame}{From keys to slots}
  \begin{definitionbox}[Hash function]
    A hash function $h$ maps every key in a universe $U$ to one of $m$ slots,
    \mbox{$h\colon U \to \{0, 1, \dots, m-1\}$}. A hash table stores key $k$ in slot $h(k)$.
  \end{definitionbox}
  \begin{examplebox}[Division method]
    With $m = 11$ and $h(k) = k \bmod 11$, the key $47$ goes to slot $3$,
    because $47 = 4 \cdot 11 + 3$.
  \end{examplebox}
  {\small\color{lcMuted}A good $h$ spreads keys evenly and is cheap to compute; choose $m$ prime
  so that patterns in the keys do not line up with the table size.\par}
\end{frame}

\begin{frame}[fragile]{Separate chaining in Python}
  \begin{columns}[T,onlytextwidth]
    \column{0.62\textwidth}
\begin{lstlisting}[language=Python]
class ChainedTable:
    def __init__(self, m=11):
        self.slots = [[] for _ in range(m)]

    def index(self, key):
        return hash(key) % len(self.slots)

    def put(self, key, value):
        chain = self.slots[self.index(key)]
        for i, (k, _) in enumerate(chain):
            if k == key:   # overwrite
                chain[i] = (key, value)
                return
        chain.append((key, value))
\end{lstlisting}
    \column{0.32\textwidth}
      \begin{itemize}\small
        \item Each slot holds a list of the pairs that hash there.
        \item \texttt{put} scans one chain, never the whole table.
        \item \texttt{get} walks the same chain and raises \texttt{KeyError} on a miss.
        \item Cost of a lookup: one hash plus the length of one chain.
      \end{itemize}
  \end{columns}
\end{frame}

% ------------------------------------------------------------------ part 2
\section{Collisions}

\begin{frame}{Linear probing, one insertion at a time}
  \framesubtitle{Insert 25, 47, 14, 36 into $m = 11$ slots with $h(k) = k \bmod 11$}
  \vspace{2mm}
  \begin{columns}[T,onlytextwidth]
    \column{0.47\textwidth}
      \begin{enumerate}\small
        \item<1-> $h(25) = 3$: slot 3 is free, store 25.
        \item<2-> $h(47) = 3$: taken, probe 4, store 47.
        \item<3-> $h(14) = 3$: probe 4, then 5, store 14.
        \item<4-> $h(36) = 3$: three probes, store in 6.
      \end{enumerate}
    \column{0.47\textwidth}
      \vspace{1mm}
      \begin{tikzpicture}[x=5.4mm,y=5.4mm]
        \foreach \i in {0,...,10} {
          \draw[lcRule,fill=white] (\i,0) rectangle ++(1,1);
          \node[font=\tiny,text=lcMuted] at (\i+0.5,-0.4) {\i};
        }
        \foreach \key/\slot/\step in {25/3/1,47/4/2,14/5/3,36/6/4} {
          \visible<\step->{
            \fill[lcAccent!14] (\slot,0) rectangle ++(1,1);
            \draw[lcAccent] (\slot,0) rectangle ++(1,1);
            \node[font=\scriptsize\bfseries,text=lcSlate] at (\slot+0.5,0.5) {\key};}
        }
        \visible<4->{\draw[lcWarning,line width=0.8pt] (3,1.35) -- (7,1.35)
          node[midway,above,font=\tiny,text=lcWarning] {a cluster of four};}
      \end{tikzpicture}
  \end{columns}
  \vspace{3mm}
  \uncover<4->{\small Keys that hash near a cluster join it: clusters grow faster the bigger they get.}
\end{frame}

\begin{frame}{How much does a collision cost?}
  \begin{theorembox}[Expected probes, uniform hashing]
    In an open-addressing table with load factor $\alpha = n/m < 1$, an unsuccessful
    search makes at most $1/(1-\alpha)$ probes in expectation.
  \end{theorembox}
  \begin{warningbox}[Deleting under open addressing]
    Do not empty a slot on delete: a later search would stop there and miss keys
    stored past it. Mark the slot with a tombstone instead.
  \end{warningbox}
\end{frame}

\begin{frame}{Chaining or open addressing?}
  \vspace{2mm}
  \centering
  \begin{tabular}{@{}l l l@{}}
    \toprule
    & \textbf{Separate chaining} & \textbf{Open addressing} \\
    \midrule
    Where colliding keys go & a list per slot & another slot of the array \\
    Load factor $\alpha$     & may exceed 1      & must stay below 1 \\
    Memory per key           & one node and pointer & none beyond the array \\
    Cache behaviour          & poor, pointers    & good, contiguous \\
    Deletion                 & remove from list  & needs tombstones \\
    \bottomrule
  \end{tabular}\par
  \vspace{4mm}
  \raggedright{\small\color{lcMuted}Python's \texttt{dict} uses open addressing; Java's
  \texttt{HashMap} uses chaining.\par}
\end{frame}

% ------------------------------------------------------------------ part 3
\section{Load factor and resizing}

\begin{frame}{Worked example: when should the table grow?}
  \begin{columns}[T,onlytextwidth]
    \column{0.5\textwidth}
      {\small A table with $m = 64$ slots holds $n = 48$ keys. How many probes does a
      failed lookup expect, and after doubling?\par}
      \vspace{1mm}
      {\small
      \begin{align*}
        \alpha &= \frac{n}{m} = \frac{48}{64} = 0.75 \\[0.3mm]
        E[\text{probes}] &\le \frac{1}{1-\alpha} = \frac{1}{0.25} = 4 \\[0.3mm]
        \alpha' &= \frac{48}{128} = 0.375 \\[0.3mm]
        E'[\text{probes}] &\le \frac{1}{0.625} = 1.6
      \end{align*}}
    \column{0.44\textwidth}
      \begin{examplebox}[Rule of thumb]
        Doubling $m$ when $\alpha$ passes $0.75$ keeps lookups under four probes.
        Each key is rehashed into the new table, so one resize costs $O(n)$.
      \end{examplebox}
      \begin{theorembox}[Amortised cost]
        With doubling, $n$ insertions cost $O(n)$ in total: $O(1)$ each, amortised.
      \end{theorembox}
  \end{columns}
\end{frame}

% ------------------------------------------------------------------ part 4
\section{Practice}

\begin{frame}{Quick question}
  \framesubtitle{Take a minute, then compare with your neighbour.}
  \vspace{1mm}
  {\small A linear-probing table has $m = 10$ and $\alpha = 0.9$. You insert one more key.
  What is the expected number of probes for the next failed lookup?\par}
  \vspace{2.5mm}
  \begin{columns}[T,onlytextwidth]
    \column{0.48\textwidth}
      \lcoption{A}{About 1, the table is nearly empty}
      \lcoption{B}{About 2}
    \column{0.48\textwidth}
      \lcoption*{C}{No bound: $\alpha = 1$, the table is full}
      \lcoption{D}{About 10}
  \end{columns}
  \vspace{1mm}
  \uncover<2->{{\small\color{lcMuted}With ten keys in ten slots there is no empty slot to stop
  a probe sequence: the bound $1/(1-\alpha)$ is undefined. Resize before this happens.\par}}
\end{frame}

\begin{frame}{Summary}
  \begin{itemize}
    \item A hash function turns a key into a slot in $O(1)$; collisions are unavoidable.
    \item Chaining keeps a list per slot; open addressing probes for another slot.
    \item The load factor $\alpha = n/m$ sets the cost: about $1/(1-\alpha)$ probes.
    \item Grow the table, usually by doubling, before $\alpha$ gets close to 1.
  \end{itemize}
  \vspace{3mm}
  \begin{definitionbox}[Key term]
    \textbf{Load factor} $\alpha$: the number of stored keys divided by the number of slots.
  \end{definitionbox}
\end{frame}

\begin{frame}{Next lecture}
  \begin{columns}[T,onlytextwidth]
    \column{0.52\textwidth}
      {\color{lcAccent}\lclabel{Lecture 8\hspace{2mm}Thursday 16 October}}\par\vspace{2mm}
      {\Large\bfseries\color{lcSlate}Binary search trees\par}
      \vspace{3mm}
      {\small Ordered keys, in-order traversal, and why a tree can
      answer questions a hash table cannot.\par}
    \column{0.42\textwidth}
      \begin{tcolorbox}[enhanced, arc=1.6mm, boxrule=0pt, colback=lcCode, left=3mm, right=3mm,
          top=2.5mm, bottom=2.5mm, boxsep=0pt, fontupper=\small]
        {\color{lcMuted}\lclabel{Before Thursday}}\par\vspace{1.5mm}
        \begin{itemize}
          \item Read chapter 12, sections 12.1 to 12.3
          \item Problem set 4 is due Friday at 17:00
          \item Office hours: Wednesday 14:00, room B214
        \end{itemize}
      \end{tcolorbox}
  \end{columns}
\end{frame}

\end{document}
